\chapter{Geometry Processing for SVG-to-Mesh Conversion}
\label{chap:geometry_conversion}

\section{Conversion role in the rendering pipeline}

In \texttt{Svg2GPU}, geometry conversion is the bridge between semantic SVG input and explicit rendering primitives.
Parsing alone is not enough for GPU pipelines.
Render backends require finite buffers (vertices, indices, style parameters), not command strings.

A compact model of the flow is:

\begin{equation}
\mathcal{S}_{svg} \xrightarrow{parse} \mathcal{R}_{typed} \xrightarrow{convert} \mathcal{G}_{gpu} \xrightarrow{draw} \mathcal{I}_{pixels},
\label{eq:conversion_pipeline}
\end{equation}

where $\mathcal{G}_{gpu}$ denotes render-ready geometry.
Correctness and performance both depend heavily on the conversion step.

\section{Tessellation fundamentals}

\subsection{Why tessellation is required}

GPU rasterizers operate on primitives, typically triangles.
Therefore continuous SVG contours must be discretized.
For fills, this means triangulation of interior regions.
For strokes, this means generating finite-width geometry around centerlines.

\subsection{Fill and stroke as distinct problems}

Although both originate from paths, fill and stroke conversion differ:

\begin{itemize}
\item fill conversion is primarily topological (inside/outside interpretation),
\item stroke conversion is primarily geometric (normals, joins, caps, width).
\end{itemize}

Treating them as separate concerns improves implementation clarity and debugging.

\section{Triangulation strategies}

\subsection{Practical selection criteria}

A triangulation strategy for browser workloads should balance:

\begin{itemize}
\item correctness on common contour types,
\item robustness under moderate edge cases,
\item manageable implementation complexity,
\item deterministic behavior.
\end{itemize}

\subsection{Delaunay and constrained variants}

Delaunay triangulation is well studied and produces high-quality meshes \cite{shewchuk_triangle}.
However, plain Delaunay does not automatically preserve polygon boundaries used by SVG fills.
Constrained variants solve this but increase algorithmic and implementation complexity.

\subsection{Ear clipping baseline}

For a thesis prototype, ear clipping is a practical baseline.
It is easier to implement and inspect, and works well for many simple polygons.
Its known limitations (degenerate contours, non-simple polygons) can be mitigated with pre-validation and controlled fallback.

\begin{table}[htbp]
\centering
\caption{Triangulation options and trade-offs}
\label{tab:triangulation_tradeoffs}
\begin{tabular}{>{\raggedright\arraybackslash}p{3.4cm} >{\raggedright\arraybackslash}p{2.2cm} >{\raggedright\arraybackslash}p{3.8cm} >{\raggedright\arraybackslash}p{4.5cm}}
\toprule
\textbf{Method} & \textbf{Complexity} & \textbf{Main advantage} & \textbf{Main limitation} \\
\midrule
Ear clipping & $O(n^2)$ & simple and transparent implementation & weaker behavior on pathological contours \\
Monotone decomposition + triangulation & $O(n \log n)$ & better asymptotic profile & higher implementation cost \\
Constrained Delaunay triangulation & $O(n \log n)$ avg. & high mesh quality with edge constraints & advanced data structure requirements \\
Triangle fan fallback & $O(n)$ & minimal overhead & invalid for concave regions \\
\bottomrule
\end{tabular}
\end{table}

\section{Numerical robustness}

\subsection{Typical failure sources}

Real SVG input may include:

\begin{itemize}
\item repeated points,
\item nearly collinear segments,
\item self-intersections,
\item tiny edges,
\item mixed winding in nested contours.
\end{itemize}

Unchecked, these issues can generate invalid triangles or unstable buffers.

\subsection{Robustness rules}

A practical robustness policy includes:

\begin{itemize}
\item finite-number checks after each conversion stage,
\item tolerance-aware geometric tests,
\item contour simplification before triangulation,
\item deterministic fallback for non-convertible cases,
\item explicit diagnostics for rejected geometry.
\end{itemize}

\subsection{Orientation support}

Signed-area orientation tests are commonly used before triangulation:

\begin{equation}
A = \frac{1}{2}\sum_{i=0}^{n-1} \left(x_i y_{i+1} - x_{i+1} y_i\right),
\label{eq:signed_area_compact}
\end{equation}

with cyclic indexing.
Orientation normalization helps preserve fill behavior consistency.

\section{Stroke mesh generation}

A stroke of width $w$ can be approximated by offsetting segment points along local normals.
If $\mathbf{t}$ is tangent direction, one normal form is:

\begin{equation}
\mathbf{n} = \frac{1}{\|\mathbf{t}\|}
\begin{bmatrix}
-t_y \\
 t_x
\end{bmatrix}.
\label{eq:stroke_normal}
\end{equation}

Offset vertices become $p_i \pm \frac{w}{2}\mathbf{n}$.
Segment strips are then stitched according to join/cap style policy.

In scope for this thesis, the focus is a robust baseline for common joins/caps, not full stylistic parity of the SVG specification.

\section{Data layout for rendering}

After conversion, geometry is packed into contiguous arrays.
Typical layouts include:

\begin{itemize}
\item position buffer (float),
\item optional index buffer,
\item style uniforms or compact per-vertex style data.
\end{itemize}

This stage affects both transfer cost and draw throughput.
Even in 2D scenes, reducing duplication and preserving stable ordering improves practical performance behavior.

\section{Failure classes and mitigation}

\begin{table}[htbp]
\centering
\caption{Failure classes during conversion}
\label{tab:conversion_failure_classes}
\begin{tabular}{>{\raggedright\arraybackslash}p{2.6cm} >{\raggedright\arraybackslash}p{4.3cm} >{\raggedright\arraybackslash}p{4.9cm}}
\toprule
\textbf{Class} & \textbf{Typical trigger} & \textbf{Mitigation} \\
\midrule
Parsing & malformed command/parameter stream & reject element and emit diagnostic \\
Topology & invalid contour structure & cleanup attempt, otherwise skip fill path \\
Numeric & NaN/overflow/underflow coordinates & normalization and finite-value checks \\
Style & unsupported style combinations & deterministic default/fallback policy \\
Render-boundary & geometry not compatible with backend assumptions & adapt layout or skip draw path \\
\bottomrule
\end{tabular}
\end{table}

This classification helps separate root causes and makes debugging more systematic.

\section{Algorithmic execution outline}

For a parsed path-like element $E$, a practical conversion routine can be summarized as:

\begin{enumerate}
\item decode command stream into geometric segments,
\item flatten curves/arcs under configured tolerance,
\item extract contours and normalize winding policy,
\item triangulate fill geometry if enabled,
\item build stroke geometry if enabled,
\item validate generated buffers before render submission.
\end{enumerate}

This order is important.
If style or topology checks are postponed too late, invalid geometry may already propagate into buffer generation.

\subsection{Deterministic fallback strategy}

When conversion fails for a contour, the system should apply deterministic fallback:
reject non-convertible fill geometry, preserve still-valid stroke paths if possible, and emit diagnostics.
This behavior is preferable to silent corruption because it keeps failure local and observable.

\section{Compact cost perspective}

Let $k$ be number of path commands and $n$ the number of vertices after flattening.
A practical preprocessing estimate is:

\begin{equation}
T_{pre} \approx T_{parse}(k) + T_{flatten}(k) + T_{triangulate}(n) + T_{stroke}(n).
\label{eq:compact_cost_model}
\end{equation}

This decomposition is useful for prioritizing optimization work.
For many practical scenes, reducing unnecessary vertex growth during flattening has a larger impact than micro-optimizing draw submission code.
Similarly, avoiding invalid contours early prevents wasted work in triangulation and buffer allocation stages.

\subsection{Tolerance configuration and vertex growth}

Flattening tolerance is one of the most sensitive parameters in the full pipeline.
Very strict tolerance improves geometric fidelity but can increase vertex count significantly, which raises both conversion and rendering cost.
Very relaxed tolerance lowers cost but can visibly distort curves.

A practical control objective is to keep approximation error bounded while limiting vertex growth:

\begin{equation}
N_{out} \leq \gamma \cdot N_{base},
\label{eq:vertex_growth_bound}
\end{equation}

where $N_{base}$ is the initial command-derived point count and $\gamma$ is a configured growth factor.
This type of guard is useful in browser workflows because it avoids accidental resource spikes on dense or malformed input.

\subsection{Preprocessing before triangulation}

Before triangulation, conversion should apply a compact preprocessing pipeline:

\begin{enumerate}
\item remove duplicate adjacent points,
\item collapse segments shorter than a minimum threshold,
\item normalize winding orientation according to fill policy,
\item split contours into independently valid rings when needed.
\end{enumerate}

These steps reduce geometric noise and improve downstream determinism.
They also decouple topological cleanup from triangulation logic, which simplifies debugging and testing.

\subsection{Diagnostics as a geometry-quality tool}

Geometry conversion quality is often assessed only by final visual output.
In practice, diagnostics are equally important.
Useful diagnostics include rejected contour count, degenerate-segment count, fallback activations, and resulting vertex/index totals.

For a thesis-scale system, these metrics serve two purposes:

\begin{itemize}
\item they expose conversion behavior under stress scenarios without requiring external profilers;
\item they make regressions visible even when visual differences are subtle.
\end{itemize}

This approach supports disciplined iteration because algorithmic changes can be evaluated against stable, measurable signals.

\section{Implications for \texttt{Svg2GPU}}

The implemented project follows these principles:

\begin{itemize}
\item parser output is typed and deterministic,
\item geometry conversion is isolated from XML parsing concerns,
\item style normalization happens before render submission,
\item conversion failures are handled explicitly,
\item architecture remains open for stronger triangulation and backend evolution.
\end{itemize}

This is aligned with practical computational geometry guidance \cite{deberg_computational_geometry} and with real-world web graphics development practice \cite{mapbox_earcut}.

\section{Implementation-oriented quality criteria}

For conversion layers used in production-like contexts, correctness must be measured with explicit quality criteria.
In \texttt{Svg2GPU}, the following criteria are considered essential:

\begin{itemize}
\item \textbf{topological consistency:} generated fill geometry should not produce unexpected overlaps for valid simple contours,
\item \textbf{numeric safety:} output buffers must contain finite values only,
\item \textbf{style coherence:} fill/stroke interpretation must match normalized style semantics,
\item \textbf{deterministic output ordering:} equal input and configuration should produce equivalent geometry order.
\end{itemize}

These criteria are intentionally practical.
They are easy to map to tests and diagnostic checks, and they align with the main risks observed in vector conversion systems.

By adopting explicit criteria, geometry conversion becomes auditable rather than heuristic.
This is important for long-term maintenance because regressions can be discussed in terms of violated properties, not only visual symptoms.

\section{Chapter summary}

This chapter established the geometry-processing strategy used by the project.
The key result is a practical and robust conversion model: triangulation/stroke generation with deterministic validation and explicit fallback behavior.
This model is sufficient for the implemented scope and provides a stable foundation for implementation and evaluation.

The next chapter maps these decisions to concrete architecture and code-level organization in \texttt{Svg2GPU}.
